\documentclass{article}

\usepackage{graphicx}
\usepackage{hyperref}
\usepackage{booktabs}
\usepackage{array}
\usepackage{geometry}
\usepackage{xcolor}
\usepackage{listings}
\usepackage{parskip}

\geometry{margin=1in}

\hypersetup{
    colorlinks=true,
    linkcolor=blue,
    urlcolor=blue,
    citecolor=blue
}

\title{\textbf{MyVCS --- Custom Version Control System}\\[0.5em]
       \large Group Project Technical Report}
\author{
    Akshat (2023B5A70994G) \and
    Kevalina Bhattacharyya (2023B4A70930G) \and
    Raymond (2023B4A70905G) \and
    Anoushka Gupta (2023B4A70903G) \and
    Shlok Upadhyay (2023B2A71166G) \and
    Moksh (2023B3A70812G)
}
\date{\today}

\begin{document}

\maketitle
\tableofcontents
\newpage

% =============================================================================
\section{Project Overview}
% =============================================================================

MyVCS is a custom Version Control System modelled after Git, built as a
multi-component system that combines a high-performance C++ core engine, a
Node.js command-line interface, a React-based web dashboard, and a Supabase
cloud backend for persistent metadata storage.  The system supports standard
version-control operations including \texttt{init}, \texttt{add},
\texttt{commit}, \texttt{status}, \texttt{diff}, \texttt{log},
\texttt{branch}, and \texttt{checkout}.

% =============================================================================
\section{Member Contributions}
% =============================================================================

The table below summarises the primary contribution of each team member as
declared in \texttt{contributions.txt}.  Any project module or task not
explicitly listed in that file has been attributed to \textbf{Raymond}.

\begin{table}[htbp]
\centering
\renewcommand{\arraystretch}{1.4}
\begin{tabular}{>{\bfseries}l l p{7cm}}
\toprule
Member & Student ID & Primary Contribution \\
\midrule
Akshat
    & 2023B5A70994G
    & Core Database \& Hashing --- SHA-1/MD5 hashing, zlib compression,
      blob storage, and the foundational \texttt{Blob} class
      (\texttt{storage.cpp}). \\
Kevalina Bhattacharyya
    & 2023B4A70930G
    & Front-End Development \& Access-Control Interface --- React dashboard
      (commits, branches, users), React Router, TailwindCSS, and
      frontend route/component-level access-control
      (\texttt{web-dashboard/}, \texttt{staging.js}). \\
Raymond
    & 2023B4A70905G
    & Supabase Integration \& Architecture --- Tree and Commit objects
      (Merkle trees, author/date metadata), Supabase cloud-backend
      integration, schema design, seed scripts, build system, and
      overall project architecture
      (\texttt{history.cpp}, \texttt{supabase-schema.sql},
       \texttt{Makefile}, \texttt{tests/}). \\
Anoushka Gupta
    & 2023B4A70903G
    & Refs, Timeline Navigation \& Security Integration ---
      \texttt{.myvcs/refs} management, branch creation, HEAD pointer,
      \texttt{log} command, automated test-account creation, auth-profile
      synchronisation, and RBAC policies for Admin/User/Viewer tiers
      (\texttt{refs.js}, \texttt{seed-auth-users.js}). \\
Shlok Upadhyay
    & 2023B2A71166G
    & Analytical Engine \& Status --- Line-by-line diff algorithm,
      file-version comparison, and the \texttt{status} command
      (Modified / Deleted / Untracked)
      (\texttt{diff.cpp}). \\
Moksh
    & 2023B3A70812G
    & CLI Interface \& Worktree Management --- Commander.js argument
      parsing, the main CLI entry point, C++ binary process-spawning,
      and the \texttt{checkout} worktree restore logic
      (\texttt{main.js}, \texttt{utils.js}). \\
\bottomrule
\end{tabular}
\caption{Member contribution breakdown}
\label{tab:contributions}
\end{table}

% =============================================================================
\section{Technical Overview}
% =============================================================================

\subsection{Architecture}

MyVCS follows a layered architecture with four distinct tiers:

\begin{itemize}
    \item \textbf{Front-End (Web Dashboard)} --- A React~18 single-page
          application built with Vite and styled using TailwindCSS.  It
          provides a dark-themed, animated interface for browsing commit
          history, branches, and user management.  Routing is handled by
          React Router~DOM and the Supabase JavaScript client is used
          directly from the browser (anon key only).

    \item \textbf{Backend / CLI} --- A Node.js~18 CLI tool
          (\texttt{myvcs}) built with Commander.js.  It orchestrates all
          VCS operations, manages the local \texttt{.myvcs/} repository
          directory, and spawns the C++ engine binaries for
          performance-critical work.  The CLI also synchronises commit
          and branch metadata to Supabase after each write operation.

    \item \textbf{C++ Core Engine} --- Three compiled binaries provide
          the computational backbone:
          \begin{itemize}
              \item \texttt{myvcs-storage} --- SHA-1 hashing and zlib
                    compression of objects.
              \item \texttt{myvcs-history} --- Merkle-tree construction
                    and commit-chain management.
              \item \texttt{myvcs-diff} --- Unified diff and working-tree
                    status comparison.
          \end{itemize}

    \item \textbf{Database (Supabase / PostgreSQL)} --- A hosted
          Supabase project stores users, commits, branches,
          access-control grants, and an audit-log transaction table.
          Row-Level Security (RLS) policies enforce role-based access at
          the database layer.
\end{itemize}

The high-level data flow is:

\begin{center}
\texttt{User $\rightarrow$ CLI (Node.js) $\rightarrow$ C++ Binaries}\\[0.3em]
\texttt{CLI $\rightarrow$ Supabase (metadata sync) $\rightarrow$ React Dashboard}
\end{center}

\subsection{User Access Control}

Access control is enforced at two levels: the Supabase Row-Level Security
layer (backend) and the React dashboard (frontend route/component guards).

Users are assigned one of three roles stored in the \texttt{users.role}
column:

\begin{table}[htbp]
\centering
\renewcommand{\arraystretch}{1.4}
\begin{tabular}{l l p{7.5cm}}
\toprule
Role & Privilege Level & Accessible Resources \\
\midrule
\textbf{admin}
    & Full access
    & Dashboard, Commits, Branches, Users table; can manage all
      \texttt{access\_control} grants and view all audit-log
      transactions. \\
\textbf{user}
    & Standard access
    & Dashboard, Commits, Branches; can insert new commits and manage
      branches; can view and update own profile only. \\
\textbf{viewer}
    & Read-only access
    & Dashboard, Commits (read-only); cannot access Branches or Users
      pages; cannot create commits or branches. \\
\bottomrule
\end{tabular}
\caption{Role-based access-control matrix}
\label{tab:rbac}
\end{table}

The helper function \texttt{can\_view\_resource(resource\_name)} is
defined in the Supabase schema and is called by every RLS policy to
enforce the matrix above uniformly across all tables.

\subsection{Module Breakdown}

\begin{table}[htbp]
\centering
\renewcommand{\arraystretch}{1.4}
\begin{tabular}{l l l p{5.5cm}}
\toprule
Module & File & Author & Description \\
\midrule
Object Storage  & \texttt{storage.cpp}   & Akshat    &
    SHA-1 hashing, zlib blob compression, \texttt{.myvcs/objects/}
    write/read. \\
Tree \& Commit  & \texttt{history.cpp}   & Raymond   &
    Merkle tree construction, commit-object serialisation, parent-chain
    linking. \\
Diff Engine     & \texttt{diff.cpp}      & Shlok     &
    Line-by-line unified diff, Modified/Deleted/Untracked status. \\
CLI Entry       & \texttt{main.js}       & Moksh     &
    Commander.js routing, C++ process-spawning, checkout worktree
    restore. \\
Staging         & \texttt{staging.js}    & Kevalina  &
    \texttt{.myvcs/index} management, \texttt{add} command, repo
    initialisation. \\
Refs Manager    & \texttt{refs.js}       & Anoushka  &
    Branch creation, HEAD pointer management, \texttt{log} traversal. \\
Supabase Client & \texttt{supabase.js}   & Raymond   &
    Auth session management, commit/branch metadata sync to cloud. \\
Web Dashboard   & \texttt{web-dashboard/} & Kevalina &
    React SPA: Dashboard, Commits, Branches, Users, Login pages with
    dark theme and animated UI. \\
DB Schema       & \texttt{supabase-schema.sql} & Raymond &
    PostgreSQL tables, RLS policies, helper functions, triggers. \\
Build System    & \texttt{Makefile}      & Raymond   &
    Root Makefile orchestrating C++ build, npm installs, tests. \\
Test Suite      & \texttt{tests/}        & Raymond   &
    Quick, basic, and stress shell-based test scripts. \\
\bottomrule
\end{tabular}
\caption{Functional module breakdown}
\label{tab:modules}
\end{table}

% =============================================================================
\section{Visuals --- System Demo Screenshots}
% =============================================================================

\begin{figure}[htbp]
    \centering
    \fbox{\rule{0pt}{5cm}\rule{12cm}{0pt}}
    \caption{System Demo --- Web Dashboard: Overview page showing commit
             count, branch count, and user count statistics cards together
             with the recent-commit feed.}
    \label{fig:demo-dashboard}
\end{figure}

\begin{figure}[htbp]
    \centering
    \fbox{\rule{0pt}{5cm}\rule{12cm}{0pt}}
    \caption{System Demo --- Commits page displaying the full commit
             history table with hash, author, message, and timestamp
             columns.}
    \label{fig:demo-commits}
\end{figure}

\begin{figure}[htbp]
    \centering
    \fbox{\rule{0pt}{5cm}\rule{12cm}{0pt}}
    \caption{System Demo --- Branches page showing branch cards with
             branch name and the associated latest commit hash.}
    \label{fig:demo-branches}
\end{figure}

\begin{figure}[htbp]
    \centering
    \fbox{\rule{0pt}{5cm}\rule{12cm}{0pt}}
    \caption{System Demo --- Users page listing team members with their
             assigned roles (admin / user / viewer) and permission
             badges.}
    \label{fig:demo-users}
\end{figure}

\begin{figure}[htbp]
    \centering
    \fbox{\rule{0pt}{5cm}\rule{12cm}{0pt}}
    \caption{System Demo --- CLI terminal session: \texttt{myvcs init},
             \texttt{myvcs add}, \texttt{myvcs commit}, \texttt{myvcs log},
             and \texttt{myvcs diff} commands executed against a sample
             repository.}
    \label{fig:demo-cli}
\end{figure}

% =============================================================================
\section{AI Integration \& API Endpoints}
% =============================================================================

\subsection{Overview}

Because MyVCS exposes its backend data through Supabase's auto-generated
REST API, any external AI agent (e.g.\ a Large Language Model tool-use
agent, a CI/CD bot, or an autonomous code-review assistant) can consume
these endpoints over HTTPS without additional server code.  Authentication
is handled via Supabase JWTs; the agent must present a valid Bearer token
obtained through the standard Supabase Auth flow.

\subsection{Authentication}

Before calling any protected endpoint the agent must authenticate:

\begin{lstlisting}[language=bash, basicstyle=\ttfamily\small,
                   breaklines=true, frame=single]
POST https://<project-ref>.supabase.co/auth/v1/token?grant_type=password
Content-Type: application/json

{
  "email": "agent@example.com",
  "password": "<password>"
}
# Response contains: access_token (JWT), token_type, expires_in
\end{lstlisting}

All subsequent requests must include:

\begin{lstlisting}[language=bash, basicstyle=\ttfamily\small,
                   breaklines=true, frame=single]
Authorization: Bearer <access_token>
apikey: <SUPABASE_ANON_KEY>
\end{lstlisting}

\subsection{Available REST Endpoints}

Supabase exposes a PostgREST-compatible REST layer at
\texttt{https://<project-ref>.supabase.co/rest/v1/}.  The endpoints
relevant to an AI agent are:

\subsubsection{Commits}

\begin{itemize}
    \item \textbf{List commits}\\
          \texttt{GET /rest/v1/commits?select=*\&order=created\_at.desc}\\
          Returns the full commit log.  An agent can use this to build
          context about the repository state before suggesting code changes.

    \item \textbf{Filter commits by author}\\
          \texttt{GET /rest/v1/commits?author\_id=eq.<author>\&select=*}\\
          Useful for per-contributor analysis or automated code-review
          assignment.

    \item \textbf{Insert a commit} (authenticated, \texttt{user} role+)\\
          \texttt{POST /rest/v1/commits}\\
          Body: \texttt{\{"commit\_hash","tree\_hash","parent\_hash",
          "author\_id","message"\}}\\
          An agent driving an automated CI pipeline can push commit
          metadata after a successful build.
\end{itemize}

\subsubsection{Branches}

\begin{itemize}
    \item \textbf{List branches}\\
          \texttt{GET /rest/v1/branches?select=*}\\
          Allows an agent to enumerate open feature branches and decide
          which to review or merge.

    \item \textbf{Update a branch pointer} (authenticated)\\
          \texttt{PATCH /rest/v1/branches?ref\_name=eq.<branch>}\\
          Body: \texttt{\{"commit\_hash": "<new-hash>"\}}\\
          An automated merge agent uses this to advance a branch ref after
          a successful merge operation.
\end{itemize}

\subsubsection{Users}

\begin{itemize}
    \item \textbf{List users} (\texttt{admin} role only)\\
          \texttt{GET /rest/v1/users?select=id,username,role,email}\\
          An orchestration agent can use this to look up user IDs when
          assigning tasks or sending notifications.
\end{itemize}

\subsubsection{Access Control}

\begin{itemize}
    \item \textbf{Query permissions for a user} (\texttt{admin} or self)\\
          \texttt{GET /rest/v1/access\_control?user\_id=eq.<uid>\&select=*}\\
          An agent enforcing least-privilege policies can read and audit
          permission grants.

    \item \textbf{Grant a permission} (\texttt{admin} only)\\
          \texttt{POST /rest/v1/access\_control}\\
          Body: \texttt{\{"user\_id","resource","permission\_type"\}}
\end{itemize}

\subsubsection{Audit Log (Transactions)}

\begin{itemize}
    \item \textbf{Retrieve audit events}\\
          \texttt{GET /rest/v1/transactions?order=created\_at.desc}\\
          An AI monitoring agent can poll this endpoint to detect
          anomalous activity patterns (e.g.\ bulk deletions, repeated
          failed auth attempts stored as transaction records).
\end{itemize}

\subsection{Example Agent Workflow}

Below is a representative task an AI agent could automate end-to-end
using the endpoints described above:

\begin{enumerate}
    \item \textbf{Authenticate} --- obtain a JWT for the agent service
          account.
    \item \textbf{Poll \texttt{/branches}} --- detect any branch whose
          \texttt{updated\_at} timestamp is more recent than the last
          check.
    \item \textbf{Fetch commits} on that branch via
          \texttt{/commits?author\_id=...} and summarise the change log
          using an LLM.
    \item \textbf{Post a review comment} (via an external notification
          webhook or email) addressed to the committer's username
          retrieved from \texttt{/users}.
    \item \textbf{Write an audit entry} to \texttt{/transactions}
          recording the review action for traceability.
\end{enumerate}

This workflow requires only the standard Supabase anon/JWT credentials and
no additional server-side code changes, making it straightforward to
integrate external AI tooling with the existing MyVCS backend.

% =============================================================================
\section{Database Schema Summary}
% =============================================================================

The Supabase PostgreSQL schema consists of five primary tables:

\begin{itemize}
    \item \textbf{users} --- Stores user accounts with
          \texttt{id} (UUID), \texttt{email}, \texttt{username}, and
          \texttt{role} (\texttt{user} $|$ \texttt{admin} $|$
          \texttt{viewer}).
    \item \textbf{commits} --- Records every commit with
          \texttt{commit\_hash}, \texttt{tree\_hash},
          \texttt{parent\_hash}, \texttt{author\_id}, and
          \texttt{message}.
    \item \textbf{branches} --- Maps branch names (\texttt{ref\_name})
          to their current \texttt{commit\_hash}.
    \item \textbf{access\_control} --- Stores per-user resource
          permission grants (\texttt{read} $|$ \texttt{write} $|$
          \texttt{admin}).
    \item \textbf{transactions} --- Append-only audit log of all
          significant actions with a \texttt{details} JSONB field.
\end{itemize}

All tables have Row-Level Security enabled.  The helper functions
\texttt{current\_user\_role()} and \texttt{can\_view\_resource()} are
registered as \texttt{SECURITY DEFINER} SQL functions to prevent privilege
escalation through direct function invocation.

% =============================================================================
\section{Build \& Deployment}
% =============================================================================

\subsection{Prerequisites}

\begin{itemize}
    \item \textbf{C++ toolchain}: \texttt{clang++} with C++17 support.
    \item \textbf{Libraries}: OpenSSL (\texttt{libssl-dev}) and zlib
          (\texttt{zlib1g-dev}).
    \item \textbf{Node.js} 18 or later.
    \item A \textbf{Supabase} project with the schema from
          \texttt{supabase-schema.sql} applied.
\end{itemize}

\subsection{Build Commands}

\begin{lstlisting}[language=bash, basicstyle=\ttfamily\small,
                   breaklines=true, frame=single]
make              # Build everything (C++ engine + npm deps)
make install      # Build + install 'myvcs' CLI globally
make test         # Run basic test suite
make test-all     # Run all tests (basic + stress)
make clean        # Remove all build artefacts
\end{lstlisting}

\end{document}
